\documentclass[11pt]{article}
\usepackage{amsmath,amssymb,amsthm}
\newtheorem{theorem}{Theorem}
\title{Exact corridor service from an inherited motion state}
\begin{document}
\maketitle

If a new duty inherits the motion of work already underway, it cannot
choose a freely placed standby. On a straight corridor the exact
replacement test uses both the minimum braking excursion and the set
of possible terminal stops.

\paragraph{Model.}
The relay stays in a closed interval \(C=[l,h]\). Its exact state
\((x,v)\) is measured at the observation cut, when a terminal region
is revealed. Its velocity is absolutely continuous, with
\(|\dot v|\le a\), \(|v|\le V\), positive \(a,V\), and deadline
\(T\ge0\) for a stop. The walls and bounds are fixed. A clearance
requirement may be encoded by shrinking \(C\) in advance. No other
agent moves in the interval, and arrival does not certify delivery.

\begin{theorem}
Write \(b=x+v|v|/(2a)\). There is no feasible path if \(x\notin C\),
\(|v|>\min(V,aT)\), or \(b\notin C\). Otherwise the exact terminal set is
\[
 C\cap\left[x+\int_0^T L(t)\,dt,\quad x+\int_0^T U(t)\,dt\right],
 \qquad
 \begin{aligned}
 L(t)&=\max(-V,v-at,-a(T-t)),\\
 U(t)&=\min(V,v+at,a(T-t)).
 \end{aligned}
\]
\end{theorem}
\begin{proof}
Speed and the Lipschitz bounds from both velocity endpoints force every
feasible velocity between \(L\) and \(U\). For \(v\ge0\), integrating
\(v-at\) from zero to \(v/a\) forces position at least \(b\) by that
time. For \(v<0\) reflect the coordinate and interchange the walls.
This proves necessity of the braking wall test.

For sufficiency assume \(v\ge0\) by reflection. The envelopes are
\(a\)-Lipschitz, within the speed bound, and have the required initial
and terminal velocities. Pure braking \(B(t)=\max(0,v-at)\) stops at
\(b\) and holds there. Since \(U\) and \(B\) are nonnegative throughout,
their velocity interpolations attain each position between \(b\) and
the clipped upper endpoint without leaving the segment between \(x\)
and the target. Since \(L=B\) up to \(v/a\) and \(L\le0=B\) thereafter,
their velocity interpolations first reach \(b\), then move monotonically
left. This attains every point between the clipped lower endpoint and
\(b\) while staying between the walls. Interpolation preserves the
kinematic bounds and endpoint velocities. This includes \(T=0\) and
degenerate endpoint ranges by using pure braking where appropriate.
\end{proof}

For arbitrary terminal-region queries at fixed \(T,C,a,V\), carry the
interval and test intersection. Distinct intervals differ on a singleton
query. This does not justify dropping velocity when the next query may
change the law or horizon, or when it needs constraint satisfaction
beyond the endpoint. With \(x=3/2\), \(a=V=1\), \(T=2\), and \(C=[0,3]\),
\(v=1\) gives \([7/4,3]\), while \(v=-1\) gives \([0,5/4]\).
At \(x=11/4,v=1\), the braking excursion instead makes the entire
interval infeasible, even for targets that lie inside the walls.

\paragraph{A stop at the observation cut can lose the mission.}
Start at rest at \(0\) inside the same corridor. At time \(2\) the
relay learns whether it must serve \(2\) or \(3\); it must be at
rest at the right point by time \(4\). With rest also required at the
observation cut, each interval grants at most one unit of displacement.
The point-\(3\) world is lost. Instead, accelerate one unit of time
and coast one unit to the common state \((3/2,1)\). Brake and hold if
the report names \(2\); coast one more unit then brake if it names
\(3\). Both controls obey the same bounds and walls. No distinguishing
observation is available before time \(2\).

\paragraph{What was checked.}
The witness rules out a comparator that forces rest at the observation
cut, not a state-space planner that retains velocity. An independent
lattice optimizer checks the extrema without walls in 1,530 parameter
cases: all envelope switches in those cases lie on its time and
velocity lattice. Wall safety follows from the continuous proof,
not the computation. This authored example is not a field test or
acceptance of R6.
\end{document}
